% =============================================================================
% Poster: GPU-Accelerated Quantum Algorithms for Two-Stage Stochastic Optimisation
% arXiv:2402.15029  |  github.com/NatLabRockies/quantum_stochastic_programming
%
% Compile with:  pdflatex poster.tex   (twice for TikZ overlays)
% Requires: texlive-science, texlive-latex-extra, texlive-fonts-extra
%           beamerposter, tcolorbox, pgfplots, tikz
% =============================================================================
\documentclass[final]{beamer}
\usepackage[size=a0,orientation=portrait,scale=1.22]{beamerposter}

% ---- packages ----------------------------------------------------------------
\usepackage{amsmath,amssymb,amsfonts}
\usepackage{graphicx}
\usepackage{booktabs}
\usepackage{tikz}
\usetikzlibrary{shapes.geometric,arrows.meta,positioning,fit,backgrounds,calc}
\usepackage{pgfplots}
\pgfplotsset{compat=1.18}
\usepackage[most]{tcolorbox}
\usepackage{array,tabularx}
\usepackage{xcolor}
\usepackage{url}
\usepackage{microtype}

% ---- colours -----------------------------------------------------------------
\definecolor{titleblue}  {RGB}{20,  40,  80}
\definecolor{accentblue} {RGB}{0,   82,  155}
\definecolor{nrelgreen}  {RGB}{0,   132, 61}
\definecolor{lightblue}  {RGB}{230, 240, 255}
\definecolor{lightgreen} {RGB}{225, 245, 230}
\definecolor{lightyellow}{RGB}{255, 252, 220}
\definecolor{warnorange} {RGB}{195, 80,  0}
\definecolor{midgray}    {RGB}{100, 100, 100}

% ---- beamer overrides --------------------------------------------------------
\setbeamertemplate{navigation symbols}{}
\setbeamertemplate{headline}{}
\setbeamertemplate{footline}{}
\setbeamercolor{background canvas}{bg=white}

% ---- tcolorbox helpers -------------------------------------------------------
\newcommand{\theoryblock}[2]{%
  \begin{tcolorbox}[
      enhanced, arc=6pt, boxrule=2pt,
      colback=lightblue, colframe=accentblue,
      colbacktitle=accentblue, coltitle=white,
      fonttitle=\bfseries\large, title={#1},
      top=6pt, bottom=6pt, left=9pt, right=9pt,
      before skip=10pt, after skip=0pt,
  ]#2\end{tcolorbox}}

\newcommand{\implblock}[2]{%
  \begin{tcolorbox}[
      enhanced, arc=6pt, boxrule=2pt,
      colback=lightgreen, colframe=nrelgreen,
      colbacktitle=nrelgreen, coltitle=white,
      fonttitle=\bfseries\large, title={#1},
      top=6pt, bottom=6pt, left=9pt, right=9pt,
      before skip=10pt, after skip=0pt,
  ]#2\end{tcolorbox}}

\newcommand{\keyblock}[2]{%
  \begin{tcolorbox}[
      enhanced, arc=6pt, boxrule=2pt,
      colback=lightyellow, colframe=warnorange,
      colbacktitle=warnorange, coltitle=white,
      fonttitle=\bfseries\large, title={#1},
      top=6pt, bottom=6pt, left=9pt, right=9pt,
      before skip=10pt, after skip=0pt,
  ]#2\end{tcolorbox}}

% ---- misc --------------------------------------------------------------------
\renewcommand{\arraystretch}{1.35}
\setlength{\tabcolsep}{7pt}

% ==============================================================================
\begin{document}
\begin{frame}[t]{}

% ==============================================================================
% HEADER
% ==============================================================================
\begin{tikzpicture}[remember picture, overlay]
  \fill[titleblue]
    ([xshift=-3cm,yshift=0.8cm]current page.north west)
    rectangle
    ([xshift=3cm,yshift=-7.4cm]current page.north east);
\end{tikzpicture}

\vspace{-0.3cm}
\begin{center}
  {\color{white}\Huge\bfseries
    GPU-Accelerated Quantum Algorithms for\\[0.3em]
    Two-Stage Stochastic Optimisation}\\[0.7cm]
  {\color{white!90!gray}\large
    C.~Rotello \quad P.~Graf \quad M.~Reynolds \quad E.\,B.~Jones
    \quad C.\,J.~Winkleblack \quad W.~Jones}\\[0.3cm]
  {\color{white!80!gray}\normalsize
    National Renewable Energy Laboratory (NREL) — NatLabRockies}\\[0.4cm]
  \begin{tikzpicture}
    \node[fill=white!20!titleblue, rounded corners=4pt, inner xsep=10pt, inner ysep=4pt]
      {\color{white}\small\textbf{arXiv:}
        \texttt{\color{yellow!80!white}2402.15029}\ [quant-ph]
        \quad\textbf{Code:}
        \texttt{\color{yellow!80!white}github.com/NatLabRockies/quantum\_stochastic\_programming}};
  \end{tikzpicture}
\end{center}
\vspace{0.55cm}

% ==============================================================================
% BODY — two equal columns
% ==============================================================================
\begin{columns}[T,totalwidth=\linewidth]

% ~~~~~~~~~~~~~~~~~~~~~~~~~~~~~~~~~~~~~~~~~~~~~~~~~~~~~~~~~~~~~~~~~~~~~~~~~~~~~~
% LEFT COLUMN — THEORY
% ~~~~~~~~~~~~~~~~~~~~~~~~~~~~~~~~~~~~~~~~~~~~~~~~~~~~~~~~~~~~~~~~~~~~~~~~~~~~~~
\begin{column}{0.492\linewidth}

% ------------------------------------------------------------------------------
\theoryblock{1 \quad Two-Stage Stochastic Unit Commitment}{

\textbf{Setting:} Decide gas generator commitment \emph{before} wind is revealed;
dispatch optimally \emph{after}.

\vspace{0.4cm}
\begin{center}
\begin{tikzpicture}[>=Stealth, node distance=2.4cm, font=\small]
  \node[draw, rounded corners=5pt, fill=accentblue!15,
        minimum width=3.2cm, minimum height=1.1cm, align=center]
        (s1) {First stage\\commit $x\!\in\!\{0,1\}^{n_g}$};
  \node[draw, rounded corners=5pt, fill=orange!25,
        minimum width=2.8cm, minimum height=1.1cm, align=center,
        right=1.8cm of s1]
        (xi) {Wind scenario\\$\xi\sim p(\xi)$};
  \node[draw, rounded corners=5pt, fill=nrelgreen!15,
        minimum width=3.2cm, minimum height=1.1cm, align=center,
        right=1.8cm of xi]
        (s2) {Second stage\\dispatch $y,\,r$};
  \draw[->] (s1) -- (xi) node[midway,above,font=\footnotesize]{decide};
  \draw[->] (xi) -- (s2) node[midway,above,font=\footnotesize]{revealed};
\end{tikzpicture}
\end{center}

\vspace{0.3cm}
\textbf{Objective} — minimise total expected cost:
\[
  \min_{x}\;\Bigl[\,\underbrace{c_g\,x}_{\text{gas cost}}
  \;+\;
  \underbrace{\mathbb{E}_{\xi}\!\left[\min_{y,r}\,\bigl\{c_w^\top y + c_r\,r\bigr\}\right]}_{\displaystyle\phi(x)\ \text{(expected value function)}}\Bigr]
\]
subject to $\displaystyle\sum_j y_j + r = d - x$,\quad $y_j \le \xi_j$,\quad $y\in\{0,1\}^{n_y}$,\quad $r\ge 0$.

\vspace{0.3cm}
\textbf{Challenge:} computing $\phi(x)$ requires integrating $Q(x,\xi)$ over all $2^{n_y}$ wind scenarios.
Classical Monte Carlo requires $O(1/\varepsilon^2)$ samples for error $\varepsilon$.
}

\vspace{0.2cm}

% ------------------------------------------------------------------------------
\theoryblock{2 \quad Discrete Quantum Annealing (DQA)}{

DQA encodes \textbf{all scenarios simultaneously} as a quantum superposition and
anneals the dispatch register to the per-scenario optimum.

\vspace{0.3cm}
\textbf{Step 1 — Encode probability distribution:}
\[
  U_\text{pdf}\,|0\rangle^{n_y} = |\psi_\xi\rangle = \sum_{\xi}\sqrt{p(\xi)}\,|\xi\rangle
\]

\textbf{Step 2 — Anneal dispatch register ($T$ steps):}
\[
  U_\text{DQA} = \prod_{t=1}^{T}\;
    \underbrace{e^{-i\gamma_t\,C(\xi)}}_{\text{phase separation}}
    \cdot
    \underbrace{e^{-i\beta_t\,B}}_{\text{mixer}}
\]
where $C(\xi)\!=\!Q(x,y,\xi)$ is the second-stage cost operator (conditioned on $|\xi\rangle$)
and $B=\sum_j X_j$ is the transverse-field mixer.

\vspace{0.3cm}
\textbf{After $T$ steps} the joint state approximates:
\[
  |\Psi\rangle \;\approx\;
  \sum_{\xi}\sqrt{p(\xi)}\;|\xi\rangle\,|y^*(\xi)\rangle,
  \qquad y^*(\xi) = \operatorname*{arg\,min}_{y}\,Q(x,y,\xi)
\]
\begin{center}
\begin{tikzpicture}
  \node[fill=accentblue!12, rounded corners=5pt, inner xsep=10pt, inner ysep=6pt]{%
    $n_y$ scenarios solved \textbf{in parallel} by a single $T$-step annealing circuit};
\end{tikzpicture}
\end{center}
}

\vspace{0.2cm}

% ------------------------------------------------------------------------------
\theoryblock{3 \quad Quantum Amplitude Estimation (QAE)}{

QAE wraps DQA in a phase-estimation circuit to extract $\phi(x)$ with
$O(1/\varepsilon)$ circuit calls — a \textbf{quadratic speedup} over classical MC.

\vspace{0.45cm}

% Circuit diagram
\begin{center}
\begin{tikzpicture}[>=Stealth, font=\small, x=1.15cm, y=0.95cm]

  % ---- QPE / QAE ancilla register ----
  \node[anchor=east] at (0,4)   {$|0\rangle^m$\;\footnotesize(QAE)};
  \draw (0.1,4) -- (10.5,4);
  % Hadamard
  \draw[fill=accentblue!25] (0.8,3.65) rectangle (1.7,4.35);
  \node at (1.25,4) {$H^{\!\otimes m}$};
  % Controlled-U boxes
  \foreach \cx/\lab in {3.2/{$\mathcal{Q}^{\!1}$}, 5.2/{$\mathcal{Q}^{\!2}$}, 7.2/{$\mathcal{Q}^{\!2^{m-1}}$}} {
    \draw[fill=accentblue!15] (\cx-0.75,3.65) rectangle (\cx+0.75,4.35);
    \node at (\cx,4) {\lab};
    \draw[dashed,->] (\cx,3.65) -- (\cx,2.35);
  }
  \node[midgray,font=\footnotesize] at (5.2,4.75) {controlled Grover iterates};
  % QFT dagger
  \draw[fill=purple!20] (8.7,3.65) rectangle (9.9,4.35);
  \node at (9.3,4) {$\mathrm{QFT}^\dagger$};
  % Measurement
  \draw (9.9,4) -- (10.5,4);
  \draw[fill=gray!30] (10.5,3.7) -- (10.5,4.3) -- (11.1,4) -- cycle;
  \node[right,font=\footnotesize] at (11.1,4) {$\hat\phi(x)$};

  % ---- PDF register ----
  \node[anchor=east] at (0,2) {$|0\rangle^{n_y}$\;\footnotesize(PDF)};
  \draw (0.1,2) -- (10.5,2);
  \draw[fill=nrelgreen!25] (0.8,1.65) rectangle (1.7,2.35);
  \node at (1.25,2) {$U_{\!\text{pdf}}$};

  % ---- Work / dispatch register ----
  \node[anchor=east] at (0,0) {$|0\rangle^{n_y}$\;\footnotesize(dispatch)};
  \draw (0.1,0) -- (10.5,0);
  \draw[fill=orange!20] (2.2,-0.4) rectangle (8.0,2.4);
  \node at (5.1,1.0) {DQA\;($T$ steps,\;$2n_y$ qubits)};

  % H on dispatch
  \draw[fill=orange!30] (1.0,-0.4) rectangle (1.9,0.4);
  \node at (1.45,0) {$H^{\!\otimes n_y}$};

  % Dots between controlled-U boxes
  \node at (6.2,4) {\Large$\cdots$};

\end{tikzpicture}
\end{center}

\vspace{0.4cm}
$\mathcal{Q} = -U_\text{DQA}\,S_0\,U_\text{DQA}^\dagger\,S_\chi$ is the Grover iterate;
measuring the QPE register gives amplitude $\sin^2\!\theta \propto \phi(x)$.

\vspace{0.3cm}
\begin{center}
\begin{tabular}{lcc}
  \toprule
  \textbf{Method} & \textbf{Evaluations} & \textbf{Precision} \\
  \midrule
  Classical Monte Carlo & $O(1/\varepsilon^2)$ & $\varepsilon$ \\
  \textbf{Quantum DQA + QAE} & $\boldsymbol{O(1/\varepsilon)}$ & $\varepsilon$ \\
  \bottomrule
\end{tabular}
\end{center}

\vspace{0.2cm}
\textbf{Qubit count:} $n = m + 2n_y$\;qubits.
At $n_y=14$ (Unit Commitment with 14 turbines): $n\!=\!28$ data qubits.
}

\end{column}

% ~~~~~~~~~~~~~~~~~~~~~~~~~~~~~~~~~~~~~~~~~~~~~~~~~~~~~~~~~~~~~~~~~~~~~~~~~~~~~~
% RIGHT COLUMN — IMPLEMENTATION
% ~~~~~~~~~~~~~~~~~~~~~~~~~~~~~~~~~~~~~~~~~~~~~~~~~~~~~~~~~~~~~~~~~~~~~~~~~~~~~~
\begin{column}{0.492\linewidth}

% ------------------------------------------------------------------------------
\implblock{4 \quad GPU Simulation Stack on NREL Kestrel}{

\begin{center}
\begin{tabular}{llp{8cm}}
  \toprule
  \textbf{Package} & \textbf{Version} & \textbf{Role} \\
  \midrule
  CUDA-Q         & 0.14.0  & Quantum kernel JIT-compilation (NVRTC) \\
  cuQuantum      & 26.3.0  & cuStateVec GPU statevector backend \\
  qiskit-aer-gpu & 0.15.1  & GPU-accelerated Qiskit Aer simulator \\
  mpi4py         & 4.1.1   & Multi-node shot distribution \\
  Cray MPICH     & 8.1.28  & GPU-aware MPI (GTL for mgpu) \\
  \bottomrule
\end{tabular}
\end{center}

\vspace{0.2cm}
\textbf{Hardware:} NREL Kestrel HPC \textemdash{} up to 8 $\times$ NVIDIA H100 SXM5 80\,GB (sm\_90)
across 2 nodes; Cray Slingshot interconnect.
}

\vspace{0.2cm}

% ------------------------------------------------------------------------------
\implblock{5 \quad Three Parallelisation Strategies}{

Noisy trajectory simulation requires \textbf{one full statevector per shot}
($2^n \times 16$\,bytes on-device). Three strategies attack this bottleneck:

\vspace{0.35cm}
\begin{center}
\begin{tabular}{p{2.8cm}p{7.2cm}p{5.5cm}}
  \toprule
  \textbf{Strategy} & \textbf{Mechanism} & \textbf{When to use} \\
  \midrule
  \textbf{mqpu}\newline(CUDA-Q) &
    Each GPU holds a full statevector copy; shots split via
    \texttt{cudaq.sample\_async(qpu\_id=i)} &
    $n_y\!\ge\!10$, single node; no MPI boilerplate \\[4pt]
  \textbf{mpi4py}\newline shot-split &
    One MPI rank per GPU (\texttt{CUDA\_VISIBLE\_DEVICES=SLURM\_LOCALID});
    ranks gather shot-count dicts &
    $n_y\!\ge\!10$, multi-node; near-ideal scaling \\[4pt]
  \textbf{mgpu}\newline(cuStateVec) &
    Statevector \emph{sharded} across GPUs via cuStateVec; enables
    $n\!>\!33$\,qubits; requires Cray GTL (\texttt{libmpi\_gtl\_cuda.so}) &
    Noiseless simulation; $n\!>\!33$\,q \\
  \bottomrule
\end{tabular}
\end{center}
}

\vspace{0.2cm}

% ------------------------------------------------------------------------------
\implblock{6 \quad Benchmark Results}{

\textbf{Noisy simulation:} Depolarising noise ($p_1\!=\!10^{-4}$, $p_2\!=\!10^{-3}$),
256 shots, 20 DQA steps, NREL Kestrel H100s.

\vspace{0.3cm}
\begin{center}
\begin{tikzpicture}
\begin{axis}[
    ybar, bar width=22pt,
    width=0.93\linewidth, height=15.5cm,
    symbolic x coords={$n_y\!=\!10$\\(20 q), $n_y\!=\!12$\\(24 q), $n_y\!=\!14$\\(28 q)},
    xtick=data,
    xticklabel style={font=\normalsize, align=center},
    x tick label as interval=false,
    ymode=log, log basis y=10,
    ylabel={Wall time (seconds)}, ylabel style={font=\normalsize},
    ymin=7, ymax=2500,
    ytick={10, 20, 50, 100, 200, 500, 1000},
    yticklabels={10, 20, 50, 100, 200, 500, 1000},
    yticklabel style={font=\normalsize},
    grid=major, grid style={gray!25, dashed},
    legend style={
        at={(0.97,0.97)}, anchor=north east,
        font=\normalsize, draw=gray!50,
        fill=white, fill opacity=0.9, text opacity=1,
        row sep=2pt,
    },
    enlarge x limits=0.22,
    clip=false,
    title={\normalsize Wall time — noisy DQA+QAE evaluation (log scale)},
    title style={font=\bfseries\normalsize},
    nodes near coords style={font=\footnotesize, /pgf/number format/fixed},
]
  % 1-GPU baseline
  \addplot[fill=gray!55, draw=gray!70] coordinates {
    ($n_y\!=\!10$\\(20 q), 46.9)
    ($n_y\!=\!12$\\(24 q), 69.2)
    ($n_y\!=\!14$\\(28 q), 808.1)
  };
  % 2-GPU mqpu
  \addplot[fill=accentblue!65, draw=accentblue!80] coordinates {
    ($n_y\!=\!10$\\(20 q), 24.9)
    ($n_y\!=\!12$\\(24 q), 36.1)
    ($n_y\!=\!14$\\(28 q), 283.2)
  };
  % 4-GPU mpi4py
  \addplot[fill=nrelgreen!75, draw=nrelgreen!90] coordinates {
    ($n_y\!=\!10$\\(20 q), 12.0)
    ($n_y\!=\!12$\\(24 q), 19.3)
    ($n_y\!=\!14$\\(28 q), 144.3)
  };
  \legend{1 GPU (baseline), 2-GPU mqpu, 4-GPU mpi4py}
\end{axis}
\end{tikzpicture}
\end{center}

\vspace{0.1cm}
\begin{center}
\begin{tabular}{ccccccc}
  \toprule
  $n_y$ & Qubits & 1-GPU (s) & 2-GPU mqpu (s) & 4-GPU mpi4py (s)
        & Su$_{2\times}$ & Su$_{4\times}$ \\
  \midrule
  10 & 20 & \phantom{0}46.9 & \phantom{0}24.9 & \phantom{0}12.0
     & 1.9$\times$ & 3.9$\times$ \\
  12 & 24 & \phantom{0}69.2 & \phantom{0}36.1 & \phantom{0}19.3
     & 1.9$\times$ & 3.6$\times$ \\
  14 & 28 & 808\phantom{.0} & 283\phantom{.0} & 144\phantom{.0}
     & 2.9$\times$ & 5.6$\times^{\dagger}$ \\
  \bottomrule
  \multicolumn{7}{l}{\footnotesize$^\dagger$ Cold-start JIT artefact;
    warm comparison $\approx$ ideal 2$\times$/4$\times$.}
\end{tabular}
\end{center}

\vspace{0.4cm}
\textbf{Noiseless mgpu} (statevector sharding, 8$\,\times\,$H100, Cray GTL):

\begin{center}
\begin{tabular}{cccc}
  \toprule
  $n_y$ & Qubits & Memory / GPU (fp32) & Wall time (s) \\
  \midrule
  16 & 32 & \phantom{0}4.3\,GB & \phantom{00}12.6 \\
  17 & 34 & \phantom{0}17.2\,GB & \phantom{00}33.5 \\
  18 & 36 & 68.7\,GB            & 122\phantom{.0} \\
  \bottomrule
\end{tabular}
\end{center}

\vspace{0.2cm}
\textbf{Memory ceiling (noisy):} fp32 statevector = $2^n \!\times\! 8$\,bytes.
Single H100 holds $\le\!33$\,q ($\approx$64\,GB). mgpu extends to 36\,q with 8 GPUs.
}

\vspace{0.2cm}

% ------------------------------------------------------------------------------
\keyblock{Key Findings}{
\begin{itemize}\setlength{\itemsep}{5pt}
  \item \textbf{mpi4py shot-splitting} achieves near-ideal $4\times$ speedup for noisy
    simulation at $n_y\!=\!10$--$12$ ($>$90\,\% parallel efficiency).
  \item \textbf{mqpu crossover} at $n_y\!=\!10$ (20\,qubits): at smaller circuits,
    async dispatch overhead exceeds compute savings; above, near-ideal $2\times$.
  \item \textbf{Practical ceiling} for noisy simulation: $n_y\!=\!14$ (28\,q) on H100;
    $n_y\!=\!16$ (32\,q, 68\,GB) causes OOM in trajectory mode.
  \item \textbf{mgpu extends noiseless simulation} to 36\,qubits ($n_y\!=\!18$)
    across 8 GPUs, but does not support noise models.
  \item \textbf{Quantum speedup:} $O(1/\varepsilon)$ QAE vs $O(1/\varepsilon^2)$
    classical Monte Carlo --- demonstrated end-to-end on NREL Kestrel H100s.
\end{itemize}
}

\end{column}
\end{columns}

% ==============================================================================
% FOOTER
% ==============================================================================
\vspace{0.5cm}
\begin{tikzpicture}[remember picture, overlay]
  \fill[titleblue]
    ([xshift=-3cm,yshift=0.5cm]current page.south west)
    rectangle
    ([xshift=3cm,yshift=3.6cm]current page.south east);
\end{tikzpicture}

\begin{columns}[T,totalwidth=\linewidth]
  \begin{column}{0.70\linewidth}
    \vspace{0.15cm}
    {\color{white}\small\raggedright
      \textbf{Reference:}\;
      C.\,Rotello, P.\,Graf, M.\,Reynolds, E.\,B.\,Jones,
      C.\,J.\,Winkleblack, W.\,Jones.
      \textit{Calculating the expected value function of a two-stage stochastic
      optimization program with a quantum algorithm.}
      arXiv:2402.15029 [quant-ph], 2024.
      \texttt{https://doi.org/10.48550/arXiv.2402.15029}

      \vspace{0.2cm}
      \textbf{Code:}\;
      \texttt{https://github.com/NatLabRockies/quantum\_stochastic\_programming}
      \quad (branch \texttt{fix/cuda-q-script})

      \vspace{0.2cm}
      \textbf{GPU implementation and benchmarks:}\;
      N.\,Sawant, NREL Kestrel HPC.\;
      CUDA-Q\,0.14.0\,/\,cuQuantum\,26.3.0\,/\,NVIDIA H100 SXM5 80\,GB.
    }
  \end{column}
  \begin{column}{0.28\linewidth}
    \vspace{0.1cm}
    {\color{white!70!titleblue}\small\raggedright
      \texttt{module load qiskit/aer-gpu}\\[2pt]
      NREL Kestrel \textemdash{} \texttt{qiskit/aer-gpu} module\\
      includes CUDA-Q, cuQuantum,\\
      Qiskit Aer GPU, mpi4py}
  \end{column}
\end{columns}

\end{frame}
\end{document}
